\section{SSH Keys}
\label{app:ssh}

SSH keys let you log in without typing a password each time, and they are how GitHub identifies
you. A key is a pair of files: a \emph{private} key that stays on your computer and must never be
shared, and a \emph{public} key (ending in \texttt{.pub}) that you give to servers.

\subsection{Create a Key Pair}

On your own computer:

\begin{lstlisting}
ssh-keygen -t ed25519 -C "your_email@uwaterloo.ca"
\end{lstlisting}

Accept the default location and choose a passphrase, which protects the private key if someone gets
the file. Show the public key with:

\begin{lstlisting}
cat ~/.ssh/id_ed25519.pub
\end{lstlisting}

To avoid typing the passphrase repeatedly, add the key to the SSH agent (many desktops do this for
you):

\begin{lstlisting}
eval "$(ssh-agent -s)"
ssh-add ~/.ssh/id_ed25519
\end{lstlisting}

\subsection{Waterloo Servers (Pelee and Rondeau)}

\begin{lstlisting}
ssh-copy-id username@pelee.math.private.uwaterloo.ca
ssh username@pelee.math.private.uwaterloo.ca
\end{lstlisting}

You should no longer be asked for a password.

\subsection{Alliance Clusters}

The Alliance recommends installing your key through its account portal, not with
\texttt{ssh-copy-id}: log in to \url{https://ccdb.alliancecan.ca}, open the SSH keys page, and
paste the content of your \emph{public} key. One upload works for all the Alliance HPC clusters.
Whether you are still asked for your second factor depends on how the key is set up; see
\url{https://docs.alliancecan.ca/wiki/SSH_Keys} and the Alliance page on multi-factor
authentication.

\subsection{GitHub}

\begin{enumerate}
\item Copy your public key (\texttt{cat \textasciitilde/.ssh/id\_ed25519.pub}).
\item On GitHub, open Settings $\rightarrow$ SSH and GPG keys $\rightarrow$ New SSH key, paste
the key, and save.
\item Test the connection:
\begin{lstlisting}
ssh -T git@github.com
\end{lstlisting}
You should see a message greeting your GitHub username.
\end{enumerate}

Clone with the SSH address (\texttt{git@github.com:user/repo.git}), not the \texttt{https}
address. For a repository that is already cloned with \texttt{https}:

\begin{lstlisting}
git remote set-url origin git@github.com:user/repo.git
git remote -v
\end{lstlisting}

On a cluster, do not copy your private key there. Create a separate key (with a passphrase) on the
login node and add its public key to GitHub too; GitHub accepts several keys. Most compute nodes
have no internet access, so run \texttt{git} commands on a login node.

\subsection{A Shortcut File}

Create \texttt{\textasciitilde/.ssh/config} on your computer:

\begin{lstlisting}
Host pelee
    HostName pelee.math.private.uwaterloo.ca
    User username

Host nibi
    HostName nibi.alliancecan.ca
    User username
    ControlMaster auto
    ControlPath ~/.ssh/cm-%C
    ControlPersist 1h
\end{lstlisting}

Now \texttt{ssh nibi} is enough. \texttt{ControlMaster} keeps the first connection open for an
hour and reuses it, so you answer the multi-factor prompt once, not for every new \texttt{ssh},
\texttt{scp} or \texttt{sshfs}.

\subsection{If It Does Not Work}

\begin{lstlisting}
chmod 700 ~/.ssh
chmod 600 ~/.ssh/id_ed25519 ~/.ssh/config
ssh -v nibi
\end{lstlisting}

The \texttt{-v} option shows which keys are tried and where the connection fails.
